\begin{description}
\item[(a)] The first step is to find an expression for the azimuth of an
  object during sunrise:

  % FIGURE NEEDED: spherical triangle with vertices Zenith, Celestial South
  % Pole and the Sun, with sides 90-a, 90+phi, 90+delta and angle 180-A at
  % the zenith; 2024 theory solutions, page 13 (problem T7).

  Using the law of cosines:
  \begin{align*}
    \cos(90^\circ + \delta)
      &= \cos(90^\circ - a)\cos(90^\circ + \phi)
       + \sin(90^\circ - a)\sin(90^\circ + \phi)\cos(180^\circ - A)\\
    \therefore -\sin(\delta)
      &= -\sin(a)\cos(\phi) - \cos(a)\cos(\phi)\cos(A)
  \end{align*}
  Since the altitude $a = 0^\circ$ at the sunrise:
  \[
    \cos(\phi) = \frac{\sin(\delta)}{\cos(A)}
  \]
  \hfill(2.0)

  Note that due to symmetry over the local meridian, the azimuth of sunrise
  corresponds to $180^\circ$ minus half of the angle given in the question
  statement (the $180^\circ$ term comes from the fact that the observer is on
  the southern hemisphere).

  The scenarios that lead to the minimum and maximum values of latitude
  correspond to the extreme values for the declination of the Sun in the
  southern hemisphere summer ($-23.5^\circ$ and $0^\circ$).

  If the declination of the Sun is $-23.5^\circ$, the latitude has to be
  \begin{align*}
    \cos(\phi) &= \frac{\sin(-23.5^\circ)}
                       {\cos\left(180^\circ - \frac{120^\circ}{2}\right)}\\
    \cos(\phi) &= -2\,\sin(-23.5^\circ)\\
    \phi &= -37.1^\circ
  \end{align*}
  \hfill(2.0)

  Only the negative solution of this equation is relevant in this case since
  the observer is in the southern hemisphere.

  If the declination of the Sun is $0^\circ$, simply plugging in this value to
  the formula will result in a latitude of $-90^\circ$. Since celestial
  objects do not rise or set at the poles, this approach would not be
  conceptually correct. However, at a point infinitesimally close to the pole,
  the Sun would still mathematically rise and set, and it would be possible to
  achieve the difference in azimuth from the problem statement with a
  declination very close to $0^\circ$. Thus, although a latitude of
  $-90^\circ$ is impossible, any latitude infinitesimally close to it would
  still mathematically be possible.

  Therefore, the range of possible latitudes for the island is the following:
  \[
    -90^\circ < \phi \leq -37.1^\circ
  \]
  \hfill(3.0)

\item[(b)] Since the latitude is constant, the following expression must be
  true:
  \[
    \cos(\phi) = \frac{\sin(\delta_0)}{\cos(A_0)}
               = \frac{\sin(\delta_{40})}{\cos(A_{40})}
  \]
  \hfill(2.0)

  It is also possible to apply the law of sines to the following triangle to
  derive another expression for the declination of the Sun on each day:

  % FIGURE NEEDED: triangle formed by the celestial equator and the ecliptic
  % meeting at the vernal point with obliquity epsilon, showing theta_0,
  % theta_40 along the ecliptic and -delta_0, -delta_40 perpendicular to the
  % equator; 2024 theory solutions, page 14 (problem T7).
  \hfill(1.0)

  \begin{align*}
    \frac{\sin(\theta_0)}{\sin(90^\circ)}
      &= \frac{\sin(-\delta_0)}{\sin(\epsilon)}\\
    \sin(\delta_0) &= -\sin(\theta_0)\sin(\epsilon)
  \end{align*}
  Analogously:
  \begin{align*}
    \frac{\sin(\theta_{40})}{\sin(90^\circ)}
      &= \frac{\sin(-\delta_{40})}{\sin(\epsilon)}\\
    \sin(\delta_{40}) &= -\sin(\theta_{40})\sin(\epsilon)
  \end{align*}
  \hfill(1.0)

  Combining the three previous expressions:
  \begin{align*}
    \sin(\delta_0) &= \frac{\sin(\delta_{40})\cos(A_0)}{\cos(A_{40})}\\
    -\sin(\theta_0)\sin(\epsilon)
      &= \frac{-\sin(\theta_{40})\sin(\epsilon)\cos(A_0)}{\cos(A_{40})}\\
    \sin(\theta_0) &= \frac{\sin(\theta_{40})\cos(A_0)}{\cos(A_{40})}\\
    \sin(\theta_0) &= k\,\sin(\theta_{40})
  \end{align*}
  \hfill(2.0)

  Where
  $k = \dfrac{\cos(A_0)}{\cos(A_{40})}
     = \dfrac{\cos(120^\circ)}{\cos(98.5^\circ)} = 3.38$.

  In order to replace one variable and solve for $\theta_0$, it is possible to
  consider that the angular velocity of the Sun on the Ecliptic is
  approximately constant due to the Earth's very low eccentricity. In that
  case, the following expression must be true:
  \begin{align*}
    \theta_{40} &= \theta_0 - \frac{40}{T_{year}} \times 360^\circ\\
    \theta_{40} &= \theta_0 - \beta
  \end{align*}
  \hfill(2.0)

  Where
  $\beta = \dfrac{40}{T_{year}} \times 360^\circ = 39.4^\circ$. Combining the
  expressions and solving for $\theta_0$:
  \begin{align*}
    \sin(\theta_0) &= k\,\sin(\theta_0 - \beta)\\
    \sin(\theta_0) &= k\,\bigl(\sin(\theta_0)\cos(\beta)
                     - \sin(\beta)\cos(\theta_0)\bigr)\\
    \sin(\theta_0)\bigl(k\,\cos(\beta) - 1\bigr)
      &= k\,\sin(\beta)\cos(\theta_0)\\
    \therefore \tan(\theta_0)
      &= \frac{k\,\sin(\beta)}{k\,\cos(\beta) - 1}
       = \frac{3.38\,\sin(39.4^\circ)}{3.38\,\cos(39.4^\circ) - 1}\\
    \theta_0 &= 53.1^\circ
  \end{align*}
  \hfill(3.0)

  Solving for $\delta_0$:
  \begin{align*}
    \sin(\delta_0) &= -\sin(\theta_0)\sin(\epsilon)\\
    \delta_0 &= -18.6^\circ
  \end{align*}
  \hfill(1.0)

  Solving for the latitude of the island:
  \begin{align*}
    \cos(\phi) &= \frac{\sin(\delta_0)}{\cos(A_0)}
      = \frac{\sin(-18.6^\circ)}
             {\cos\left(180^\circ - \frac{120^\circ}{2}\right)}\\
    \phi &= -50.4^\circ
  \end{align*}
  \hfill(1.0)

  Note that only the negative solution to the equation above is relevant since
  it is known that the island is in the southern hemisphere.
\end{description}
